%% CD01
\begin{frame}[t]{Comprendre une petite variation}
\coursanalyse
Une fonction de plusieurs variables reçoit plusieurs informations à la fois.
\begin{cours}{La question du chapitre}
Quand l'entrée change un peu, comment prévoir la variation de la sortie ?
\end{cours}
\begin{center}\begin{tikzpicture}
\node at(0,0){$\vec a\longmapsto f(\vec a)$};\node at(7,0){$\vec a+\vec h\longmapsto f(\vec a+\vec h)$};
\draw[-{Stealth},thick,IMTCyan](2,0)--(4.6,0);\node[above]at(3.3,.3){petit déplacement};
\end{tikzpicture}\end{center}
\vfill\textbf{Fil conducteur :} remplacer localement la variation par une application linéaire.
\note{Cours original, pages PDF 49–55. Résultats admis ; expliciter le sens des formules sans démonstration.}
\end{frame}

%% CD02
\begin{frame}[t]{Entrée, sortie et composantes}
\coursanalyse
\begin{notion}{Fonction vectorielle}
Pour $U\subset\Rn n$, une fonction $f:U\to\Rn p$ s'écrit
\[\vec x=\vcol{x_1\\\vdots\\x_n},\qquad f(\vec x)=\vcol{f_1(\vec x)\\\vdots\\f_p(\vec x)}.\]
Les $f_i$ sont les \textbf{fonctions composantes} de $f$.
\end{notion}
\begin{resultat}{Deux dimensions à distinguer}
$n$ coordonnées à l'entrée ; $p$ coordonnées à la sortie. Si $p=1$, la fonction est à valeurs réelles.
\end{resultat}
\note{Cours original, pages PDF 12–14. Résultats admis ; expliciter le sens des formules sans démonstration.}
\end{frame}

%% CD03
\begin{frame}[t]{Commencer par le domaine}
\coursanalyse
\begin{cours}{Avant de calculer}
Repérer les valeurs pour lesquelles chaque expression est définie.
\end{cours}
\begin{exemple}{Deux contraintes usuelles}
$f(x,y)=\ln(3-x^2-y^2)$ : domaine $x^2+y^2<3$.

\medskip
$g(x,y)=\dfrac1{x-y}$ : domaine $x\ne y$.
\end{exemple}
\vfill Un dénominateur doit rester non nul ; l'argument d'un logarithme doit être strictement positif.
\note{Cours original, pages PDF 12–14,115–117. Résultats admis ; expliciter le sens des formules sans démonstration.}
\end{frame}

%% CD04
\begin{frame}[t]{Mesurer la taille d’un déplacement}
\coursanalyse
\begin{notion}{Norme et distance euclidiennes}
Pour $\vec h\in\Rn n$,
\[\norm{\vec h}=\sqrt{h_1^2+\cdots+h_n^2}.\]
La distance est $D(\vec x,\vec a)=\norm{\vec x-\vec a}$.
\end{notion}
\begin{resultat}{Le sens de « petit »}
$\vec h\to\zero$ signifie que \textbf{toutes} ses coordonnées tendent vers zéro.
\end{resultat}
\vfill La norme est un réel ; le déplacement est un vecteur.
\note{Cours original, pages PDF 3–11. Résultats admis ; expliciter le sens des formules sans démonstration.}
\end{frame}

%% CD05
\begin{frame}[t]{Une boule décrit la proximité}
\coursanalyse
\begin{notion}{Centre $\vec a$ et rayon $r>0$}
\textbf{Boule ouverte :} $B(\vec a,r)=\{\vec x:\norm{\vec x-\vec a}<r\}$.

\medskip
\textbf{Boule fermée :} $\overline B(\vec a,r)=\{\vec x:\norm{\vec x-\vec a}\le r\}$.
\end{notion}
\begin{center}\begin{tikzpicture}[scale=.65]
\fill[IMTCyan!10](0,0)circle(1.5);\draw[dashed,IMTCyan,thick](0,0)circle(1.5);\fill(0,0)circle(2pt)node[below]{$\vec a$};\draw(0,0)--(1.5,0)node[midway,above]{$r$};\node at(0,-2){bord exclu};
\fill[IMTGreen!18](7,0)circle(1.5);\draw[IMTNavy,thick](7,0)circle(1.5);\fill(7,0)circle(2pt)node[below]{$\vec a$};\draw(7,0)--(8.5,0)node[midway,above]{$r$};\node at(7,-2){bord inclus};
\end{tikzpicture}\end{center}
\note{Cours original, pages PDF 3–5. Résultats admis ; expliciter le sens des formules sans démonstration.}
\end{frame}

%% CD06
\begin{frame}[t]{Un ouvert laisse de la place autour de chaque point}
\coursanalyse
\begin{notion}{Partie ouverte}
$U\subset\Rn n$ est ouvert si
\[\forall\vec a\in U,\quad\exists r>0,\quad B(\vec a,r)\subset U.\]
\end{notion}
\begin{center}\begin{tikzpicture}[scale=.7]
\fill[IMTCyan!10](-2,-1.5)rectangle(2,1.5);\draw[dashed,IMTCyan,thick](-2,-1.5)rectangle(2,1.5);\draw[IMTOrange,thick](.8,.2)circle(.8);\fill(.8,.2)circle(2pt)node[below]{$\vec a$};\node at(-1.3,1){$U$};
\node[anchor=west,align=left]at(4,.2){Un petit déplacement\\dans toute direction\\reste dans $U$.};
\end{tikzpicture}\end{center}
\vfill Le rayon peut dépendre du point choisi.
\note{Cours original, pages PDF 6–7. Résultats admis ; expliciter le sens des formules sans démonstration.}
\end{frame}

%% CD07
\begin{frame}[t]{Voisinage et point intérieur}
\coursanalyse
\begin{notion}{Voisinage de $\vec a$}
$V$ est un voisinage de $\vec a$ s'il contient une boule ouverte centrée en $\vec a$.
\[\exists r>0,\qquad B(\vec a,r)\subset V.\]
\end{notion}
\begin{resultat}{Deux formulations utiles}
$\vec a$ est \textbf{intérieur} à $U$ si $U$ est un voisinage de $\vec a$.

\medskip
$U$ est \textbf{ouvert} si chacun de ses points est intérieur.
\end{resultat}
\vfill Un voisinage n'est pas nécessairement un ouvert.
\note{Cours original, pages PDF 8–9. Résultats admis ; expliciter le sens des formules sans démonstration.}
\end{frame}

%% CD08
\begin{frame}[t]{Une limite doit convenir à toutes les approches}
\coursanalyse
Soit $\vec a$ un point d'accumulation de $U$ : on peut l'approcher par des points de $U$ distincts de lui.
\begin{notion}{Limite $\vec\ell$ en $\vec a$}
Pour tout $\varepsilon>0$, il existe $\eta>0$ tel que, pour tout $\vec x\in U$,
\[0<\norm{\vec x-\vec a}<\eta\Longrightarrow\norm{f(\vec x)-\vec\ell}<\varepsilon.\]
On note $\displaystyle\lim_{\vec x\to\vec a}f(\vec x)=\vec\ell$.
\end{notion}
\begin{theoreme}{Composante par composante}
Une fonction vectorielle converge si et seulement si toutes ses composantes convergent ; leurs limites forment la colonne limite.
\end{theoreme}
\note{Cours original, pages PDF 15–16,30–31. Résultats admis ; expliciter le sens des formules sans démonstration.}
\end{frame}

%% CD09
\begin{frame}[t]{Deux chemins peuvent exclure une limite}
\coursanalyse
\begin{exemple}{Une fonction définie hors de l'origine}
\[q(x,y)=\frac{xy}{x^2+y^2}.\]
\textbf{Sur l'axe horizontal :} $q(t,0)=0$ pour $t\ne0$.

\medskip
\textbf{Sur la diagonale :} $q(t,t)=\dfrac12$ pour $t\ne0$.
\end{exemple}
\begin{resultat}{Conclusion}
Les deux valeurs limites diffèrent : $q$ n'a pas de limite à l'origine.
\end{resultat}
\vfill Plusieurs chemins de même limite ne suffisent pas, à eux seuls, à prouver une limite.
\note{Cours original, pages PDF 35–37. Résultats admis ; expliciter le sens des formules sans démonstration.}
\end{frame}

%% CD10
\begin{frame}[t]{La continuité : la limite rejoint la valeur}
\coursanalyse
\begin{notion}{Continuité en $\vec a\in U$}
Pour tout $\varepsilon>0$, il existe $\eta>0$ tel que, pour tout $\vec x\in U$,
\[\norm{\vec x-\vec a}<\eta\Longrightarrow\norm{f(\vec x)-f(\vec a)}<\varepsilon.\]
\end{notion}
Si $\vec a$ est un point d'accumulation, cela signifie
\[\lim_{\vec x\to\vec a}f(\vec x)=f(\vec a).\]
\begin{resultat}{Pour une fonction vectorielle}
$f$ est continue si et seulement si chacune de ses composantes est continue.
\end{resultat}
\note{Cours original, pages PDF 32–33,37–38. Résultats admis ; expliciter le sens des formules sans démonstration.}
\end{frame}

%% CD11
\begin{frame}[t]{Reconnaître les fonctions continues usuelles}
\coursanalyse
\begin{theoreme}{Règles de continuité}
\textbf{Somme et produit :} des fonctions réelles continues restent continues.

\medskip
\textbf{Quotient :} il reste continu là où le dénominateur ne s'annule pas.

\medskip
\textbf{Composition :} $g\circ f$ est continue en $\vec a$ si $f$ l'est en $\vec a$ et $g$ en $f(\vec a)$.
\end{theoreme}
\vfill Les polynômes sont continus ; les fonctions usuelles se composent sur leur domaine de définition.
\note{Cours original, pages PDF 34–44. Résultats admis ; expliciter le sens des formules sans démonstration.}
\end{frame}

%% CD12
\begin{frame}[t]{Un reste négligeable au premier ordre}
\coursanalyse
\begin{notion}{Petit $o$ de la norme}
\[r(\vec h)=o(\norm{\vec h})\Longleftrightarrow
\lim_{\vec h\to\zero,\,\vec h\ne\zero}\frac{\norm{r(\vec h)}}{\norm{\vec h}}=0.\]
Le reste devient petit \textbf{par rapport à la taille du déplacement}.
\end{notion}
\begin{exemple}{Un reste quadratique}
$r(\vec h)=\norm{\vec h}^2$ est négligeable devant $\norm{\vec h}$.
\end{exemple}
\vfill Le dénominateur est une norme, donc un scalaire : on ne divise jamais par un vecteur.
\note{Cours original, pages PDF 44. Résultats admis ; expliciter le sens des formules sans démonstration.}
\end{frame}

%% CD13
\begin{frame}[t]{En une variable : la tangente approche la courbe}
\coursanalyse
Pour une fonction dérivable, $f(a+h)\simeq f(a)+f'(a)h$.
\begin{center}\begin{tikzpicture}[x=1.1cm,y=.48cm]
\draw[axes](-.4,0)--(3.6,0)node[right]{$x$};\draw[axes](0,-.3)--(0,6.5)node[above]{$f(x)$};
\draw[courb,domain=0:2.5,samples=60]plot(\x,{\x*\x});\draw[IMTOrange,thick,domain=.5:3]plot(\x,{2*\x-1});
\fill(1,1)circle(2pt)node[above left]{$A$};\draw[pointille](1,0)node[below]{$a$}--(1,1);
\node[anchor=west,text=IMTNavy]at(4,5.8){courbe $y=x^2$};\node[anchor=west,text=IMTOrange]at(4,3.8){tangente en $a=1$};
\end{tikzpicture}\end{center}
\vfill En plusieurs variables, la partie linéaire remplace la multiplication par $f'(a)$.
\note{Cours original, pages PDF 47–49,59–60. Résultats admis ; expliciter le sens des formules sans démonstration.}
\end{frame}

%% CD14
\begin{frame}[t]{La différentiabilité : une approximation linéaire}
\coursanalyse
Dans la suite, $U$ est un ouvert de $\Rn n$, $\vec a\in U$ et $f:U\to\Rn p$.
\begin{notion}{Différentiabilité en $\vec a$}
Il existe une application \textbf{linéaire} $\Lin L:\Rn n\to\Rn p$ telle que
\[f(\vec a+\vec h)=f(\vec a)+\Lin L(\vec h)+o(\norm{\vec h})\]
lorsque $\vec h\to\zero$.
\end{notion}
\begin{resultat}{Lire les trois termes}
Valeur de départ \quad + \quad variation linéaire \quad + \quad reste négligeable.
\end{resultat}
\vfill Pour $\vec h$ assez petit, $\vec a+\vec h$ reste dans $U$.
\note{Cours original, pages PDF 49–50. Résultats admis ; expliciter le sens des formules sans démonstration.}
\end{frame}

%% CD15
\begin{frame}[t]{La différentielle est une application}
\coursanalyse
\begin{theoreme}{Unicité de la partie linéaire}
Si $f$ est différentiable en $\vec a$, l'application linéaire précédente est unique.
\end{theoreme}
\begin{notion}{Notation du cours}
Cette application est la \textbf{différentielle} $\df f{\vec a}$.
\[\df f{\vec a}:\Rn n\to\Rn p,\qquad\vec h\longmapsto\df f{\vec a}(\vec h).\]
\end{notion}
\vfill $\vec a$ est le point où l'on se place ; $\vec h$ est le déplacement auquel on applique la différentielle.
\note{Cours original, pages PDF 49–53. Résultats admis ; expliciter le sens des formules sans démonstration.}
\end{frame}

%% CD16
\begin{frame}[t]{Le modèle affine local}
\coursanalyse
\begin{resultat}{Prévoir la valeur voisine}
\[f(\vec a+\vec h)\simeq f(\vec a)+\df f{\vec a}(\vec h).\]
\end{resultat}
\begin{cours}{Deux objets à distinguer}
\textbf{Différentielle :} $\vec h\mapsto\df f{\vec a}(\vec h)$ est linéaire.

\medskip
\textbf{Approximation de $f$ :} $\vec x\mapsto f(\vec a)+\df f{\vec a}(\vec x-\vec a)$ est affine.
\end{cours}
\vfill Le signe $\simeq$ indique une approximation ; l'égalité exacte contient un reste.
\note{Cours original, pages PDF 49–54. Résultats admis ; expliciter le sens des formules sans démonstration.}
\end{frame}

%% CD17
\begin{frame}[t]{Un exemple d’approximation au premier ordre}
\coursanalyse
\begin{exemple}{$f(x,y)=x^2+3y^2$, au point $\vec a=\vcol{1\\1}$}
Pour un déplacement $\vec h=\vcol{h\\k}$,
\[f(1+h,1+k)=4+2h+6k+\underbrace{h^2+3k^2}_{\text{reste quadratique}}.\]
\end{exemple}
\begin{resultat}{Le modèle à conserver}
\[\df f{\vec a}\!\left(\vcol{h\\k}\right)=2h+6k,\qquad f(1+h,1+k)\simeq4+2h+6k.\]
\end{resultat}
\note{Cours original, pages PDF 49–55. Résultats admis ; expliciter le sens des formules sans démonstration.}
\end{frame}

%% CD18
\begin{frame}[t]{Constante, linéaire, affine : les cas immédiats}
\coursanalyse
\begin{theoreme}{Différentielles de référence}
\textbf{Constante :} si $f(\vec x)=\vec b$, alors $\df f{\vec a}=\Lin 0$.

\medskip
\textbf{Linéaire :} si $f=\Lin L$, alors $\df f{\vec a}=\Lin L$.

\medskip
\textbf{Affine :} si $f(\vec x)=\Lin L(\vec x)+\vec b$, alors $\df f{\vec a}=\Lin L$.
\end{theoreme}
\vfill Pour une fonction affine, le modèle au premier ordre est exact : le reste est nul.
\note{Cours original, pages PDF 54–57. Résultats admis ; expliciter le sens des formules sans démonstration.}
\end{frame}

%% CD19
\begin{frame}[t]{Différentiable implique continue}
\coursanalyse
\begin{theoreme}{Une implication à retenir}
$f$ différentiable en $\vec a$ \quad$\Longrightarrow$\quad $f$ continue en $\vec a$.
\end{theoreme}
\begin{exemple}{La réciproque est fausse}
$f(\vec x)=\norm{\vec x}$ est continue à l'origine, mais n'y est pas différentiable.

\medskip
Sur la droite $t\mapsto t\vec e_1$, elle devient $t\mapsto|t|$ : il y a un angle à l'origine.
\end{exemple}
\vfill Être proche de la valeur de départ ne suffit pas à avoir un modèle linéaire.
\note{Cours original, pages PDF 57–60. Résultats admis ; expliciter le sens des formules sans démonstration.}
\end{frame}

%% CD20
\begin{frame}[t]{Une dérivée partielle : une seule variable change}
\coursanalyse
\begin{notion}{Dérivée partielle selon $x_j$}
\[\frac{\partial f}{\partial x_j}(\vec a)=\lim_{t\to0}\frac{f(\vec a+t\vec e_j)-f(\vec a)}{t},\]
si cette limite existe. Les autres coordonnées restent fixes.
\end{notion}
\begin{exemple}{$f(x,y)=x^2+3y^2$}
\[\frac{\partial f}{\partial x}(x,y)=2x,\qquad\frac{\partial f}{\partial y}(x,y)=6y.\]
\end{exemple}
\vfill On dérive par rapport à une variable ; les autres jouent le rôle de constantes.
\note{Cours original, pages PDF 76–78. Résultats admis ; expliciter le sens des formules sans démonstration.}
\end{frame}

%% CD21
\begin{frame}[t]{Dériver une fonction vectorielle}
\coursanalyse
\begin{theoreme}{Composante par composante}
Si $f=\vcol{f_1\\\vdots\\f_p}$, alors, lorsque les dérivées existent,
\[\frac{\partial f}{\partial x_j}=\vcol{\partial f_1/\partial x_j\\\vdots\\\partial f_p/\partial x_j}.\]
\end{theoreme}
\begin{resultat}{La taille ne change pas à la sortie}
Une dérivée partielle d'une fonction à $p$ composantes est un vecteur colonne à $p$ composantes.
\end{resultat}
\note{Cours original, pages PDF 77–79,83–85. Résultats admis ; expliciter le sens des formules sans démonstration.}
\end{frame}

%% CD22
\begin{frame}[t]{Dériver suivant un vecteur donné}
\coursanalyse
\begin{notion}{Dérivée suivant $\vec u$}
\[D_{\vec u}f(\vec a)=\lim_{t\to0}\frac{f(\vec a+t\vec u)-f(\vec a)}{t},\]
si cette limite existe.
\end{notion}
\begin{cours}{Le rôle du vecteur $\vec u$}
Il fixe une \textbf{direction} et une \textbf{échelle de déplacement}.

\medskip
Il n'est pas nécessairement unitaire.
\end{cours}
\vfill Pour comparer des directions à vitesse égale, on choisit des vecteurs unitaires.
\note{Cours original, pages PDF 75–77. Résultats admis ; expliciter le sens des formules sans démonstration.}
\end{frame}

%% CD23
\begin{frame}[t]{Une différentielle, toutes les directions}
\coursanalyse
\begin{theoreme}{Si $f$ est différentiable en $\vec a$}
Pour tout vecteur $\vec u$,
\[D_{\vec u}f(\vec a)=\df f{\vec a}(\vec u).\]
En particulier,
\[\frac{\partial f}{\partial x_j}(\vec a)=\df f{\vec a}(\vec e_j).\]
\end{theoreme}
\vfill Une seule application linéaire rassemble les variations suivant toutes les directions.
\note{Cours original, pages PDF 79–82. Résultats admis ; expliciter le sens des formules sans démonstration.}
\end{frame}

%% CD24
\begin{frame}[t]{Des dérivées partielles ne suffisent pas}
\coursanalyse
\begin{exemple}{Compléter la fonction de l'exemple des deux chemins}
Posons $q(0,0)=0$ et $q(x,y)=xy/(x^2+y^2)$ ailleurs.

\medskip
Sur chaque axe, $q$ est nulle : ses deux dérivées partielles à l'origine valent $0$.

\medskip
Sur la diagonale, $q(t,t)=1/2$ pour $t\ne0$ : $q$ n'est pas continue à l'origine.
\end{exemple}
\begin{resultat}{Conclusion}
$q$ n'est pas différentiable à l'origine, malgré l'existence de ses deux dérivées partielles.
\end{resultat}
\note{Cours original, pages PDF 79–81. Résultats admis ; expliciter le sens des formules sans démonstration.}
\end{frame}

%% CD25
\begin{frame}[t]{La jacobienne représente la différentielle}
\coursanalyse
\begin{notion}{Dans les bases canoniques}
Pour $f:U\subset\Rn n\to\Rn p$ différentiable en $\vec a$,
\[J_f(\vec a)=\begin{pmatrix}
\dfrac{\partial f_1}{\partial x_1}(\vec a)&\cdots&\dfrac{\partial f_1}{\partial x_n}(\vec a)\\
\vdots&&\vdots\\
\dfrac{\partial f_p}{\partial x_1}(\vec a)&\cdots&\dfrac{\partial f_p}{\partial x_n}(\vec a)
\end{pmatrix}.\]
\end{notion}
\vfill\textbf{$p$ lignes :} les composantes de sortie.

\medskip
\textbf{$n$ colonnes :} les variables d'entrée.
\note{Cours original, pages PDF 85–86. Résultats admis ; expliciter le sens des formules sans démonstration.}
\end{frame}

%% CD26
\begin{frame}[t]{Lire les lignes et les colonnes}
\coursanalyse
\begin{exemple}{$F(x,y)=\vcol{x^2+y\\xy-y^2}$}
\[J_F(x,y)=\begin{pmatrix}2x&1\\y&x-2y\end{pmatrix}.\]
\end{exemple}
\begin{cours}{Deux lectures}
\textbf{Ligne 1 :} les dérivées de la première composante $x^2+y$.

\medskip
\textbf{Colonne 2 :} la dérivée vectorielle selon $y$, soit $\vcol{1\\x-2y}$.
\end{cours}
\note{Cours original, pages PDF 83–86. Résultats admis ; expliciter le sens des formules sans démonstration.}
\end{frame}

%% CD27
\begin{frame}[t]{Appliquer la différentielle à un déplacement}
\coursanalyse
\begin{theoreme}{Dans les bases canoniques}
\[\coord{\df f{\vec a}(\vec h)}{\mathbf e'}=J_f(\vec a)\coord{\vec h}{\mathbf e}.\]
\end{theoreme}
\begin{exemple}{Pour $F(x,y)=\vcol{x^2+y\\xy-y^2}$ en $\vec a=\vcol{1\\1}$}
\[\df F{\vec a}\!\left(\vcol{h\\k}\right)
=\begin{pmatrix}2&1\\1&-1\end{pmatrix}\vcol{h\\k}=\vcol{2h+k\\h-k}.\]
\end{exemple}
\note{Cours original, pages PDF 85–86. Résultats admis ; expliciter le sens des formules sans démonstration.}
\end{frame}

%% CD28
\begin{frame}[t]{Ne pas confondre les objets}
\coursanalyse
\begin{cours}{La fonction}
$f$ associe une sortie à chaque point de son domaine ; elle peut être non linéaire.
\end{cours}
\begin{cours}{La différentielle en un point}
$\df f{\vec a}$ est une application \textbf{linéaire} qui agit sur les déplacements.
\end{cours}
\begin{cours}{La jacobienne en ce point}
$J_f(\vec a)$ est la \textbf{matrice} de cette application dans les bases canoniques.
\end{cours}
\vfill Fixer le point $\vec a$ avant de faire agir la différentielle.
\note{Cours original, pages PDF 85–86. Résultats admis ; expliciter le sens des formules sans démonstration.}
\end{frame}

%% CD29
\begin{frame}[t]{Justifier la différentiabilité}
\coursanalyse
\begin{cours}{Pour une expression usuelle}
\textbf{1.} Vérifier le domaine et se placer sur un ouvert.

\medskip
\textbf{2.} Repérer les briques différentiables : constantes, applications affines et fonctions usuelles d'une variable.

\medskip
\textbf{3.} Appliquer les règles de somme, produit, quotient et composition, présentées dans la dernière partie.
\end{cours}
\begin{resultat}{L'ordre à conserver}
Justifier la différentiabilité, puis utiliser la jacobienne pour calculer l'action de la différentielle.
\end{resultat}
\note{Cours original, pages PDF 54–57,105–117. Résultats admis ; expliciter le sens des formules sans démonstration.}
\end{frame}

%% CD30
\begin{frame}[t]{À valeurs réelles : une forme linéaire}
\coursanalyse
Si $f:U\subset\Rn n\to\R$ est différentiable en $\vec a$,
\begin{theoreme}{Expression de la différentielle}
\[\df f{\vec a}(\vec h)=\sum_{j=1}^n\frac{\partial f}{\partial x_j}(\vec a)h_j.\]
\end{theoreme}
\begin{exemple}{$f(x,y)=x^2+3y^2$, en $\vec a=\vcol{1\\1}$}
\[\df f{\vec a}\!\left(\vcol{h\\k}\right)=2h+6k.\]
\end{exemple}

\note{Cours original, pages PDF 87–88. Résultats admis ; expliciter le sens des formules sans démonstration.}
\end{frame}

%% CD31
\begin{frame}[t]{Le gradient : représenter par un vecteur}
\coursanalyse
\begin{notion}{Gradient}
Pour $f:U\to\R$ différentiable en $\vec a$, le gradient $\grad f(\vec a)$ est l'unique vecteur tel que
\[\forall\vec h\in\Rn n,\qquad\df f{\vec a}(\vec h)=\ps{\grad f(\vec a)}{\vec h}.\]
\end{notion}
\begin{resultat}{Deux objets liés, mais distincts}
La différentielle est une \textbf{application linéaire}.

\medskip
Le gradient est un \textbf{vecteur} qui la représente grâce au produit scalaire.
\end{resultat}
\note{Cours original, pages PDF 88–90. Résultats admis ; expliciter le sens des formules sans démonstration.}
\end{frame}

%% CD32
\begin{frame}[t]{Gradient colonne et jacobienne ligne}
\coursanalyse
\begin{notion}{Dans la base canonique}
\[\grad f(\vec a)=\vcol{\dfrac{\partial f}{\partial x_1}(\vec a)\\\vdots\\\dfrac{\partial f}{\partial x_n}(\vec a)}.\]
\end{notion}
\begin{resultat}{La jacobienne est une matrice, pas un vecteur}
Pour une fonction scalaire, $J_f(\vec a)$ a une ligne et $n$ colonnes :
\[J_f(\vec a)=\bigl(\grad f(\vec a)\bigr)^\top.\]
\end{resultat}
\note{Cours original, pages PDF 88–90. Résultats admis ; expliciter le sens des formules sans démonstration.}
\end{frame}

%% CD33
\begin{frame}[t]{Voir le gradient sur des courbes de niveau}
\coursanalyse
Une courbe de niveau rassemble les points où $f$ prend la même valeur.
\begin{center}\begin{tikzpicture}[x=.75cm,y=.75cm]
\draw[axes](-2.5,0)--(3,0)node[right]{$x$};\draw[axes](0,-1.7)--(0,2.4)node[above]{$y$};
\draw[IMTCyan,thick](0,0)ellipse[x radius=2,y radius=1.1547];\draw[IMTCyan!45](0,0)ellipse[x radius=1.4142,y radius=.8165];
\fill(1,1)circle(2pt)node[below right]{$\vec a$};\draw[-{Stealth},IMTOrange,very thick](1,1)--(1.4,2.2)node[above right]{$\grad f(\vec a)$};\draw[IMTNavy,thick](-.2,1.4)--(2.2,.6);
\node[anchor=west,align=left]at(4,1){$f(x,y)=x^2+3y^2$\\[2mm]$\vec a=\vcol{1\\1}$, $\grad f(\vec a)=\vcol{2\\6}$};
\end{tikzpicture}\end{center}
\vfill Le gradient est normal à la tangente au niveau en ce point. La flèche est à échelle réduite.
\note{Cours original, pages PDF 89–90. Résultats admis ; expliciter le sens des formules sans démonstration.}
\end{frame}

%% CD34
\begin{frame}[t]{La direction de croissance maximale}
\coursanalyse
\begin{theoreme}{Parmi les vecteurs unitaires}
Si $f:U\to\R$ est différentiable en $\vec a$,
\[D_{\vec u}f(\vec a)=\ps{\grad f(\vec a)}{\vec u},\qquad\norm{\vec u}=1.\]
Si $\grad f(\vec a)\ne\zero$, le maximum vaut $\norm{\grad f(\vec a)}$, atteint pour
\[\vec u=\frac{\grad f(\vec a)}{\norm{\grad f(\vec a)}}.\]
\end{theoreme}
\begin{resultat}{Si le gradient est nul}
Toutes les dérivées directionnelles sont nulles. Cela ne signifie pas que la fonction est constante.
\end{resultat}
\note{Cours original, pages PDF 89–90. Résultats admis ; expliciter le sens des formules sans démonstration.}
\end{frame}

%% CD35
\begin{frame}[t]{Vecteur donné ou direction unitaire ?}
\coursanalyse
\begin{exemple}{$f(x,y)=x^2+3y^2$, en $\vec a=\vcol{1\\1}$}
Ici $\grad f(\vec a)=\vcol{2\\6}$ et $\vec v=\vcol{3\\4}$ a pour norme $5$.
\end{exemple}
\begin{resultat}{Conserver la convention de la question}
\textbf{Suivant $\vec v$ :} $D_{\vec v}f(\vec a)=2\times3+6\times4=30$.

\medskip
\textbf{Suivant le vecteur unitaire $\vec v/5$ :} $D_{\vec v/5}f(\vec a)=6$.
\end{resultat}
\vfill Normaliser change la valeur : annoncer le vecteur effectivement utilisé.
\note{Cours original, pages PDF 76,89–90. Résultats admis ; expliciter le sens des formules sans démonstration.}
\end{frame}

%% CD36
\begin{frame}[t]{Un résultat de référence : la norme au carré}
\coursanalyse
\begin{theoreme}{Pour $s(\vec x)=\norm{\vec x}^2$}
$s$ est différentiable sur $\Rn n$ et
\[\df s{\vec a}(\vec h)=2\ps{\vec a}{\vec h}.\]
Son gradient est $\grad s(\vec a)=2\vec a$.
\end{theoreme}
\begin{resultat}{Distinguer norme et norme au carré}
La norme au carré est différentiable même à l'origine.

\medskip
La norme, elle, n'y est pas différentiable.
\end{resultat}
\note{Cours original, pages PDF 108–109. Résultats admis ; expliciter le sens des formules sans démonstration.}
\end{frame}

%% CD37
\begin{frame}[t]{Sommes et multiples}
\coursanalyse
\begin{theoreme}{Combinaison linéaire}
Si $f,g:U\to\Rn p$ sont différentiables en $\vec a$ et $\lambda,\mu\in\R$, alors
\[\df{(\lambda f+\mu g)}{\vec a}=\lambda\df f{\vec a}+\mu\df g{\vec a}.\]
\end{theoreme}
\begin{cours}{Même règle sur les matrices}
\[J_{\lambda f+\mu g}(\vec a)=\lambda J_f(\vec a)+\mu J_g(\vec a).\]
\end{cours}
\vfill On additionne des fonctions ayant le même espace d'arrivée.
\note{Cours original, pages PDF 105–106. Résultats admis ; expliciter le sens des formules sans démonstration.}
\end{frame}

%% CD38
\begin{frame}[t]{Produit de fonctions réelles}
\coursanalyse
\begin{theoreme}{Règle du produit}
Si $u,v:U\to\R$ sont différentiables en $\vec a$,
\[\df{(uv)}{\vec a}=v(\vec a)\df u{\vec a}+u(\vec a)\df v{\vec a}.\]
\end{theoreme}
\begin{exemple}{La fonction $w(x,y)=x\,y^2$}
\[\df w{\vec a}\!\left(\vcol{h\\k}\right)=a_2^2h+2a_1a_2k.\]
\end{exemple}
\vfill Les valeurs $u(\vec a)$ et $v(\vec a)$ sont des scalaires ; les différentielles restent des applications.
\note{Cours original, pages PDF 107–108. Résultats admis ; expliciter le sens des formules sans démonstration.}
\end{frame}

%% CD39
\begin{frame}[t]{La règle de la chaîne}
\coursanalyse
Soient $U\subset\Rn n$ et $V\subset\Rn p$ ouverts, $f:U\to V$ et $g:V\to\Rn q$.
\begin{theoreme}{Différencier une composition}
Si $f$ est différentiable en $\vec a$ et $g$ en $f(\vec a)$, la composée est différentiable en $\vec a$ et
\[\df{(g\circ f)}{\vec a}=\df g{f(\vec a)}\circ\df f{\vec a}.\]
\end{theoreme}
\begin{resultat}{L'ordre et le point d'évaluation}
D'abord la différentielle de $f$ en $\vec a$.

\medskip
Ensuite celle de $g$ en \textbf{$f(\vec a)$}, et non en $\vec a$.
\end{resultat}
\note{Cours original, pages PDF 110–112. Résultats admis ; expliciter le sens des formules sans démonstration.}
\end{frame}

%% CD40
\begin{frame}[t]{Voir les deux transformations successives}
\coursanalyse
\begin{center}\begin{tikzpicture}
\node at(0,0){$\vec h\in\Rn n$};\node at(5,0){$\vec k\in\Rn p$};\node at(10,0){$\vec\ell\in\Rn q$};
\draw[-{Stealth},IMTCyan,thick](1.3,0)--(3.7,0)node[midway,above]{$\df f{\vec a}$};
\draw[-{Stealth},IMTOrange,thick](6.3,0)--(8.7,0)node[midway,above]{$\df g{f(\vec a)}$};
\end{tikzpicture}\end{center}
\begin{theoreme}{Version matricielle}
\[J_{g\circ f}(\vec a)=J_g(f(\vec a))\,J_f(\vec a).\]
\end{theoreme}
\begin{resultat}{Contrôler les dimensions}
$(q\times p)\,(p\times n)=q\times n$.

\medskip
Les matrices se multiplient ; les applications se composent avec $\circ$.
\end{resultat}
\note{Cours original, pages PDF 112–114. Résultats admis ; expliciter le sens des formules sans démonstration.}
\end{frame}

%% CD41
\begin{frame}[t]{Une composition, deux jacobiennes}
\coursanalyse
\begin{exemple}{$f(x,y)=\vcol{x+2y\\x-y}$ et $g(u,v)=u^2+2v$}
\[J_f=\begin{pmatrix}1&2\\1&-1\end{pmatrix},\qquad J_g(u,v)=\begin{pmatrix}2u&2\end{pmatrix}.\]
Pour $\vec a=\vcol{1\\0}$, on a $f(\vec a)=\vcol{1\\1}$.
\end{exemple}
\begin{resultat}{Appliquer la chaîne au bon point}
\[J_{g\circ f}(\vec a)=\begin{pmatrix}2&2\end{pmatrix}\begin{pmatrix}1&2\\1&-1\end{pmatrix}=\begin{pmatrix}4&2\end{pmatrix}.\]
\end{resultat}
\note{Cours original, pages PDF 110–114. Résultats admis ; expliciter le sens des formules sans démonstration.}
\end{frame}

%% CD42
\begin{frame}[t]{Le long d’une trajectoire}
\coursanalyse
\begin{theoreme}{Dériver une fonction observée sur un chemin}
Si $\vec\gamma$ est dérivable en $t$ et $f$ différentiable en $\vec\gamma(t)$,
\[\frac{d}{dt}(f\circ\vec\gamma)(t)=\df f{\vec\gamma(t)}\bigl(\vec\gamma'(t)\bigr).\]
Si $f$ est réelle, cela vaut $\ps{\grad f(\vec\gamma(t))}{\vec\gamma'(t)}$.
\end{theoreme}
\begin{exemple}{$f(x,y)=x^2+3y^2$, $\vec\gamma(t)=\vcol{t\\1-t}$}
La vitesse vaut $\vec\gamma'(t)=\vcol{1\\-1}$ ; la dérivée de $f\circ\vec\gamma$ vaut $8t-6$.
\end{exemple}
\note{Cours original, pages PDF 112–114. Résultats admis ; expliciter le sens des formules sans démonstration.}
\end{frame}

%% CD43
\begin{frame}[t]{Inverse et quotient : le dénominateur compte}
\coursanalyse
\begin{theoreme}{Fonctions réelles différentiables en $\vec a$}
\textbf{Inverse, si $v(\vec a)\ne0$ :}
\[\df{(1/v)}{\vec a}=-\frac{1}{v(\vec a)^2}\df v{\vec a}.\]
\textbf{Quotient, si $v(\vec a)\ne0$ :}
\[\df{(u/v)}{\vec a}=\frac{v(\vec a)\df u{\vec a}-u(\vec a)\df v{\vec a}}{v(\vec a)^2}.\]
\end{theoreme}
\vfill Le dénominateur non nul en $\vec a$ reste non nul dans un voisinage assez petit.
\note{Cours original, pages PDF 115–117. Résultats admis ; expliciter le sens des formules sans démonstration.}
\end{frame}

%% CD44
\begin{frame}[t]{Composer avec une fonction usuelle}
\coursanalyse
Pour $u:U\to\R$ différentiable en $\vec a$, la chaîne donne :
\begin{cours}{Une règle par ligne}
\textbf{Exponentielle :}\quad $\df{(\exp\circ u)}{\vec a}=e^{u(\vec a)}\df u{\vec a}$.

\bigskip
\textbf{Logarithme, $u(\vec a)>0$ :}\quad $\df{(\ln\circ u)}{\vec a}=\dfrac{\df u{\vec a}}{u(\vec a)}$.

\bigskip
\textbf{Racine, $u(\vec a)>0$ :}\quad $\df{(\sqrt{\,\cdot\,}\circ u)}{\vec a}=\dfrac{\df u{\vec a}}{2\sqrt{u(\vec a)}}$.
\end{cours}
\vfill Les conditions assurent une formule valable dans un voisinage du point.

\note{Cours original, pages PDF 110–117. Résultats admis ; expliciter le sens des formules sans démonstration.}
\end{frame}

%% CD45
\begin{frame}[t]{Une démarche de calcul en quatre étapes}
\coursanalyse
\begin{cours}{Du domaine au modèle local}
\textbf{1. Domaine :} où la fonction est-elle définie ? Le point est-il intérieur ?

\medskip
\textbf{2. Différentiabilité :} quelles règles s'appliquent et sous quelles conditions ?

\medskip
\textbf{3. Calcul :} dérivées partielles, puis jacobienne ; gradient si la sortie est réelle.

\medskip
\textbf{4. Action :} fixer $\vec a$, appliquer à $\vec h$, puis écrire l'approximation affine.
\end{cours}
\vfill Contrôler les dimensions et garder les coordonnées en colonnes.
\note{Cours original, pages PDF 49–117. Résultats admis ; expliciter le sens des formules sans démonstration.}
\end{frame}

%% CD46
\begin{frame}[t]{Estimer une variation sans tout recalculer}
\coursanalyse
\begin{exemple}{$f(x,y)=x^2+3y^2$, près de $\vec a=\vcol{1\\1}$}
Prenons $\vec h=\vcol{0{,}02\\-0{,}01}$.
\[\Delta f\simeq\df f{\vec a}(\vec h)=2(0{,}02)+6(-0{,}01)=-0{,}02.\]
Donc $f(1{,}02;0{,}99)\simeq3{,}98$.
\end{exemple}
\begin{resultat}{Interprétation}
Le modèle prévoit une petite baisse. La valeur exacte est $3{,}9807$ ; l'écart $0{,}0007$ vient des termes quadratiques.
\end{resultat}
\note{Cours original, pages PDF 49–55. Résultats admis ; expliciter le sens des formules sans démonstration.}
\end{frame}

%% CD47
\begin{frame}[t]{Relier les trois représentations}
\coursanalyse
Pour une fonction \textbf{différentiable}, au même point $\vec a$ :
\begin{cours}{Application intrinsèque}
$\df f{\vec a}$ agit linéairement sur les déplacements.
\end{cours}
\begin{cours}{Matrice dans les bases canoniques}
$J_f(\vec a)$ calcule cette action sur les colonnes de coordonnées.
\end{cours}
\begin{resultat}{Vecteur si la sortie est réelle}
$\grad f(\vec a)$ représente la différentielle par le produit scalaire.

\medskip
$D_{\vec u}f(\vec a)=\df f{\vec a}(\vec u)=\ps{\grad f(\vec a)}{\vec u}$.
\end{resultat}
\note{Cours original, pages PDF 85–90. Résultats admis ; expliciter le sens des formules sans démonstration.}
\end{frame}

%% CD48
\begin{frame}[t]{Les points de vigilance}
\coursanalyse
\begin{cours}{Avant de conclure}
\textbf{Des partielles existent :} cela ne garantit pas la différentiabilité.

\medskip
\textbf{Une direction est donnée :} ne pas la normaliser sans l'annoncer.

\medskip
\textbf{Une composition apparaît :} respecter l'ordre et évaluer en $f(\vec a)$.

\medskip
\textbf{Un quotient apparaît :} vérifier que le dénominateur reste non nul.
\end{cours}
\vfill Le cours fournit les notions et les outils ; leur mise en œuvre reste le travail du TD.
\note{Cours original, pages PDF 49–117. Résultats admis ; expliciter le sens des formules sans démonstration.}
\end{frame}

%% CD49
\begin{frame}[t]{Complément — Fermé, borné, compact}
\coursanalyse
\begin{notion}{Trois notions distinctes}
\textbf{Fermé :} le complémentaire dans $\Rn n$ est ouvert.

\medskip
\textbf{Borné :} la partie est contenue dans une boule de rayon fini.

\medskip
\textbf{Compact dans $\Rn n$ :} la partie est fermée et bornée.
\end{notion}
\begin{exemple}{Deux boules de rayon $r>0$}
$B(\vec a,r)$ est bornée, mais n'est pas fermée.

$\overline B(\vec a,r)$ est fermée et bornée, donc compacte.
\end{exemple}
\note{Cours original, pages PDF 7,10–11. Résultats admis ; expliciter le sens des formules sans démonstration.}
\end{frame}

%% CD50
\begin{frame}[t]{Complément — Dériver dans une autre base}
\coursanalyse
\begin{notion}{Base $\mathbf b=(\vec b_1,\ldots,\vec b_n)$}
La dérivée partielle relativement à cette base est
\[\partial_{\mathbf b,j}f(\vec a)=D_{\vec b_j}f(\vec a).\]
\end{notion}
\begin{theoreme}{Si $f$ est différentiable en $\vec a$}
Pour $\vec h=\sum_j h_j\vec b_j$,
\[\df f{\vec a}(\vec h)=\sum_j h_j\partial_{\mathbf b,j}f(\vec a).\]
\end{theoreme}
\vfill Les dérivées partielles usuelles correspondent à la base canonique.
\note{Cours original, pages PDF 77–83. Résultats admis ; expliciter le sens des formules sans démonstration.}
\end{frame}

%% CD51
\begin{frame}[t]{Complément — La matrice dépend des bases}
\coursanalyse
\begin{notion}{Matrice de la différentielle}
Pour une base $\mathbf b$ au départ et $\mathbf c$ à l'arrivée, on note
\[M=\Lin{\mathfrak M}\bigl(\df f{\vec a};\mathbf b,\mathbf c\bigr).\]
\end{notion}
\begin{theoreme}{Action sur les coordonnées}
\[\coord{\df f{\vec a}(\vec h)}{\mathbf c}=M\coord{\vec h}{\mathbf b}.\]
\end{theoreme}
\vfill La différentielle ne change pas avec les bases ; sa représentation matricielle change. La jacobienne correspond aux bases canoniques.
\note{Cours original, pages PDF 83–86. Résultats admis ; expliciter le sens des formules sans démonstration.}
\end{frame}

%% CD52
\begin{frame}[t]{Complément — Applications bilinéaires}
\coursanalyse
\begin{theoreme}{Différentielle d'une application bilinéaire}
Une application bilinéaire $\Lin B:\Rn n\times\Rn m\to\Rn p$ est différentiable et
\[\df{\Lin B}{(\vec a,\vec b)}(\vec h,\vec k)=\Lin B(\vec h,\vec b)+\Lin B(\vec a,\vec k).\]
\end{theoreme}
\begin{resultat}{Le principe}
Faire varier un argument à la fois, puis additionner les variations linéaires.
\end{resultat}
\vfill Le produit scalaire est un exemple d'application bilinéaire réelle.
\note{Cours original, pages PDF 106–108. Résultats admis ; expliciter le sens des formules sans démonstration.}
\end{frame}
